\begin{thebibliography}{99}
%\label{sec:bibliographie}

\vfill

\bibitem{fortran} J.~W.~Backus, R.J.~Beeber, S.~Best, R.~Goldberg, 
L.~M.~Haibt, H.~L.~Herick, R.~A.~Nelson, D.~Sayre, P.~B.~Sheridan, H.~Stern,
I.~Ziller, R.~A.~Hugues \& R.~Nutt.
{\em The Fortran automatic coding system}, 
Proceedings Western Joint Computer Conference, Volume 11, 
pages 188--198, 1957.

\bibitem{mandelbrot1} M.~Barnsley, R.~Devaney, B.~Mandelbrot,
H.-O.~Peitgen, D.~Saupe \& R.~Voss.
{\em The science of fractal images}, Springer-Verlag, Berlin, 312 pages, 1988.

\bibitem{bertho} C. Bertho.
{\em Télégraphes et téléphones, de Valmy au microprocesseur}, 
Le Livre de Poche, Paris, 541 pages, 1981.

\bibitem{benoit1} P. Benoit.
``Calcul, algèbre et marchandise'', \rm in \cite{serres1}, pages 197--221, 1989.

\bibitem{bohm} C. Böhm.
``Macchina calcolatrice digitale a programma con programma preordinato fisso con
tastiera algebrica ridotta atta a comporre formule mediante la
combinazione dei singoli elementi simbolici'', Patent application No.~13567, Milan,
October 1$^{\rm st}$, 1952, 26 pages.

\bibitem{boutet} M. Boutet de Monvel.
{\em Chansons de France pour les petits enfants}, 
Gautier-Languereau, Paris, 77 pages, 1981.

\thispagestyle{empty}

\bibitem{terminologie-info} Commission ministérielle de terminologie de l'informatique.
{\em Glossaire des termes officiels de l'informatique}, AFNOR, Tour Europe, 92080 Paris
La Défense Cedex 7, 21 pages, 1989 (2$^{\rm nd}$ edition).

\bibitem{conchon} S. Conchon \& J.-C. Filliâtre.
{\em Apprendre à programmer avec OCaml. Algorithmes et structures de données}, 
Eyrolles, Paris, 2014.

\bibitem{CousotIA} P. Cousot \& R. Cousot.
``Abstract interpretation: a unified lattice model for static analysis
of programs by construction or approximation of fixpoints'', Conference
Record of the 4$^{\rm th}$ ACM Symposium on Principles of Programming
Languages, Los Angeles, California, U.S.A., p. 238--252, 1977.

\bibitem{LT} P. Cousot.
{\em Introduction à l'algorithmique numérique et à la 
programmation en Pascal}, McGraw-Hill, Paris, 621 pages, 1988.

%\bibitem{LX} P. Cousot.
%{\em Introduction à l'algorithmique et à la 
%programmation en Pascal, leçons 1 à 6}, Cours de l'École Polytechnique, 
%Palaiseau, 322 pages, 1988.

\bibitem{escher2} H.~S.~M.~Coxeter, M.~Emmer, R.~Penrose \& M.~L.~Teuber.
{\em M.~C.~Escher: Art and Science}, North-Holland, Amsterdam, 
402 pages, 1986.

\bibitem{cundy} H.~M.~Cundy \& A.~R.~Rollett.
{\em Mathematical models}, Tarquin publications (Stradbroke, Diss, Norfolk, G.B.), 
286 pages, 1981 (3$^{\rm rd}$ edition).

\bibitem{solfege} A.~Danhauser.
{\em Théorie de la musique (édition revue et corrigée par H. Rabaud)}, 
Éditions Henry Lemoine, Paris, 128  pages, 1929.

\bibitem{golomb1} S.~W.~Golomb.
``A geometric proof of a famous identity'', 
{\em Mathematical gazette}, Vol. 69, p. 198--200, 1965.

\bibitem{grunbaum} B.~Grünbaum \& G.~C.~Shephard.
{\em Tilings and patterns}, W.~H.~Freeman \& Co., New York, 700 pages, 1987.

\bibitem{guitel} G.~Guitel.
{\em Histoire comparée des numérations écrites}, 
Flammarion, Paris, 857 pages, 1975.

\bibitem{horemis} S.~Horemis.
{\em Optical and geometrical patterns and designs}, Dover, New York, 1970.

\bibitem{iverson1} K.~E.~Iverson.
{\em A Programming Language}, J.~Wiley \& Sons, 1962.

\bibitem{wirth2} K.~Jensen \& N.~Wirth.
{\em Pascal, user manual and report}, Springer-Verlag (Heidelberg), 
167 pages, 1976 (2$^{\rm nd}$ edition).

\bibitem{kernighan1} B. W. Kernighan \& D. M. Ritchie.
{\em The C programming language}, Prentice-Hall (Englewood Cliffs, New Jersey, U.S.A.),
228 pages, 1978.

\bibitem{knuth2} D. E. Knuth.
{\em The art of computer programming, Volume 1 / Fundamental algorithms}, 
Addison-Wesley (Reading, Massachusetts, USA), 
634 pages, 1973 (2$^{\rm nd}$ edition).

\bibitem{knuth1} D. E. Knuth.
{\em The \TeX book}, Addison-Wesley (Reading, Massachusetts, USA), 
483 pages, 1987 (12$^{\rm th}$ edition).

\bibitem{lamport1} L. Lamport.
{\em \LaTeX\ a document preparation system}, Addison-Wesley (Reading, Massachusetts, USA), 
242 pages, 1986.

\bibitem{ocaml} X. Leroy, D. Doligez, A. Frisch, J. Garrigue, D. Rémy, K. Sivaramakrishnan
\& J. Vouillon. {\em The OCaml system: documentation and user's manual}, Inria, 
\url{https://ocaml.org/manual/}.

\bibitem{levy1} P. Lévy.
``L'invention de l'ordinateur'', \rm in \cite{serres1}, pages 515--535, 1989.

\bibitem{escher1} J. L. Locher.
{\em The world of M. C. Escher}, Abrams (New York, USA), 1971\/\/;
{\footnotesize French translation:} J. L. Locher,
{\em Le monde de M. C. Escher}, Éditions du Chêne (Paris), 
293 pages, 1986 (12$^{\rm th}$ edition).

\bibitem{mandelbrot2} B. Mandelbrot.
{\em Les objets fractals} {\footnotesize followed by} {\em Survol du langage
fractal}, Flammarion, Nouvelle collection scientifique (Paris), 1989
(3$^{\rm rd}$ edition).

\bibitem{lisp} J.~McCarthy, P.~W.~Abrahams, D.~J.~Edwards, T.~P.~Hart
\& M.~I.~Levin.
{\em LISP 1.5 Programmer's manual}, 
MIT Press, 1965.

\bibitem{algol60} P. Naur {\em et al.}  (Eds.).
``Revised report on the algorithmic language Algol 60'', 
{\em Communications of the Association for Computing Machinery}, 
volume 6, pages 1--17, 1963.

\bibitem{odds1} F. C. Odds.
``Spirolaterals'', {\em Mathematics teacher}, volume 66, pages 121--124, 1973.

\bibitem{logo} S. Papert.
{\em Mindstorms: children, computers and powerful ideas}, Harvester Press, 
1980.

\bibitem{python} Python Software Foundation.
{\em The Python Language Reference}, official documentation of Python~3,
\url{https://docs.python.org/3/reference/}.

\bibitem{python-tutoriel} Python Software Foundation.
{\em Le tutoriel Python}, official documentation of Python~3 (French translation),
\url{https://docs.python.org/fr/3/tutorial/}.

\bibitem{ritter1} J. Ritter.
``Babylone -1800'', \rm in \cite{serres1}, pages 17--37, 1989.

\bibitem{rutishauser} H. Rutishauser.
``Automatische Rechenplanfertigung bei programmgesteuerten
Rechenmaschinen'', {\em Mitteilungen aus dem Inst. für angew. Math. an der ETH 
Zürich}, No. 3, Birkhäuser, Basel, 1952, 45 pages.

\bibitem{sabatier1} R. Sabatier.
{\em Le livre des chansons de France}, Gallimard, Paris, 157 pages, 1984.

\bibitem{sabatier2} R. Sabatier.
{\em Deuxième livre des chansons de France et d'ailleurs}, 
Gallimard, Paris, 165 pages, 1986.

\bibitem{sabatier3} C. Sabatier \& R. Sabatier.
{\em Troisième livre des chansons de France}, 
Gallimard, Paris, 165 pages, 1987.

\bibitem{sedgewick1} R. Sedgewick.
{\em Algorithms}, Addison-Wesley, Reading, Massachusetts, U.S.A., 
650 pages, 1988 (2$^{\rm nd}$ edition), {\footnotesize translated into French by
J.-M. Moreau:} R. Sedgewick, {\em Algorithmes}, InterÉditions, Paris, 1990.

\bibitem{serres1} M. Serres.
{\em Éléments d'histoire des sciences},  Bordas, Paris, 576 pages, 1989.

\bibitem{thevenet1} S. Thévenet, A. Garioud \& N. Pitot.
{\em Mathématiques CM2},  Bordas, Paris, 160 pages, 1986.

\bibitem{voderberg} H.~Voderberg.
``Zur Zerlegung der Umgebung eines ebenen Bereiches in kongruente'', 
{\em Jahresber. Deutsch. Math.-Verein.}, 
volume 46, pages 229--231, 1936;
``Zur Zerlegung der Ebene in kongruente Bereiche in Form einer Spirale'', 
{Ibid.}, 
volume 47, pages 159--160, 1937

%\bibitem{vuillemin1} J. Vuillemin.
%Algorithmes, nombres, machines, langages, {\em École Polytechnique},
%251 pages, 1990.

\bibitem{weis} P. Weis \& X. Leroy.
{\em Le langage Caml}, InterÉditions, Paris, 1993; 2$^{\rm nd}$ edition, Dunod, 1999.

\bibitem{wilkes1} M. V. Wilkes, D. J. Wheeler \& S. Gill. 
``The preparation of programs for an electronic digital computer: with special reference to 
the EDSAC and the use of a library of subroutines'', Addison-Wesley, Cambridge, 
Massachusetts, 170 pages, 1951.

\bibitem{wilson} L.~B.~Wilson \& R.~G.~Clark.
{\em Comparative programming languages},  Addison-Wesley publishing company,
 Reading, Massachusetts, 379 pages, 1988.

\bibitem{wirth1} N. Wirth.
``The programming language Pascal'', {\em Acta Informatica}, Vol. 1, Springer-Verlag (Heidelberg), 
pages 35--63, 1971.

\end{thebibliography}

